\documentclass[11pt,a4paper]{article}

\usepackage{amsmath, amssymb, amsthm}
\usepackage{physics}
\usepackage{bm}
\usepackage{siunitx}
\usepackage{booktabs}
\usepackage{array}
\usepackage[margin=2cm]{geometry}
\usepackage{enumitem}
\usepackage{hyperref}

\renewcommand{\arraystretch}{1.6}

\title{Technical Note: Multi-Sphere T-Matrix Theory\\[4pt]
\large Mathematical Formulation and Computational Implementation}
\author{N.\ Moteki}
\date{Last updated: 2026-03-26 (corresponds to MSTMforCAS.jl v0.3.0)}

\begin{document}
\maketitle

\begin{abstract}
This document provides a self-contained description of the mathematical theory behind the Multi-Sphere T-Matrix (MSTM) method as implemented in MSTMforCAS.jl.
It covers the VSWF expansion formalism, Lorenz--Mie theory for single spheres, the multi-sphere interaction equation, VSWF translation addition theorem, the Complex BiConjugate Gradient (CBICG) iterative solver, the FFT-accelerated translation algorithm, and the output quantities (scattering amplitudes and cross-section efficiencies).
The implementation follows the formulation of Mackowski (1991, 1994, 1996), Xu (1996), Mackowski \& Mishchenko (2011), and Mackowski \& Kolokolova (2022).
Scattering amplitude conventions follow Bohren \& Huffman (1983, hereafter BH83).
\end{abstract}

\tableofcontents

%% ============================================================
\section{Problem Setup}
\label{sec:problem-setup}

\subsection{Geometry and incident field}

Consider $N_S$ homogeneous, non-overlapping spheres embedded in an infinite, lossless, homogeneous medium of real refractive index $n_\mathrm{med}$.
Sphere $i$ ($i = 1, \dots, N_S$) has centre position $\bm{r}_i$, radius $a_i$, and complex refractive index $m_\mathrm{sphere}$ (common to all spheres in the parameter-sweep mode of this code).
A time-harmonic plane wave ($e^{-i\omega t}$ convention, following BH83) propagates along the $+z$ axis in the medium with vacuum wavelength $\lambda_0$ and wavenumber in the medium
\begin{equation}
  k = \frac{2\pi \, n_\mathrm{med}}{\lambda_0}.
  \label{eq:wavenumber}
\end{equation}

A positive imaginary part of $m_\mathrm{sphere}$ corresponds to absorption.
All lengths are non-dimensionalised by $k$: positions become $k\bm{r}_i$, radii become size parameters $x_i = k a_i$, and the relative refractive index entering all scattering formulas is $m_\mathrm{rel} = m_\mathrm{sphere} / n_\mathrm{med}$.
Tildes are suppressed hereafter; all equations are in dimensionless units unless stated otherwise.

The volume-equivalent sphere radius $a_\mathrm{eff}$ and its size parameter $x_\mathrm{eff}$ are defined as
\begin{equation}
  a_\mathrm{eff} = \left(\sum_{i=1}^{N_S} a_i^3\right)^{\!1/3}, \qquad
  x_\mathrm{eff} = k\,a_\mathrm{eff}.
  \label{eq:a-eff}
\end{equation}
These appear in the normalisation of cross-section efficiencies (Sec.~\ref{sec:output}).

\subsection{Vector spherical wave function expansion}

The electromagnetic field in each region is expanded in vector spherical wave functions (VSWFs) $\bm{M}_{mn}^{(\nu)}$, $\bm{N}_{mn}^{(\nu)}$ of degree $m$ and order $n$ ($n = 1, 2, \dots$; $m = -n, \dots, n$), where $\nu = 1$ denotes regular (Bessel $j_n$) waves and $\nu = 3$ denotes outgoing (Hankel $h_n^{(1)}$) waves.
The normalisation convention for these VSWFs follows that of Mishchenko et al.\ (2002) and Mackowski (1994).
The scattered field from sphere $i$ is expanded in outgoing VSWFs centred at $\bm{r}_i$:
\begin{equation}
  \bm{E}_\mathrm{sca}^{(i)}(\bm{r}) = \sum_{n=1}^{N_i} \sum_{m=-n}^{n}
  \left[ a_{mn}^{(i)}\,\bm{M}_{mn}^{(3)}(\bm{r} - \bm{r}_i)
       + b_{mn}^{(i)}\,\bm{N}_{mn}^{(3)}(\bm{r} - \bm{r}_i) \right]
  \label{eq:sca-expansion}
\end{equation}
where $N_i$ is the truncation order for sphere $i$.
The field exciting sphere $i$ --- the incident plane wave plus the scattered fields from all other spheres --- is likewise expanded in regular VSWFs $\bm{M}_{mn}^{(1)}, \bm{N}_{mn}^{(1)}$ centred at $\bm{r}_i$.

\subsection{Left--right circular (lr\_tran) basis}

Following Mackowski (1996, \S 2), the code works in the left--right circular VSWF basis obtained by the linear combinations
\begin{equation}
  \tilde{\bm{N}}_{mn,L}^{(\nu)} = \bm{N}_{mn,1}^{(\nu)} + \bm{N}_{mn,2}^{(\nu)}, \qquad
  \tilde{\bm{N}}_{mn,R}^{(\nu)} = \bm{N}_{mn,1}^{(\nu)} - \bm{N}_{mn,2}^{(\nu)}
  \label{eq:lr-basis}
\end{equation}
where the subscripts 1, 2 refer to the TE and TM modes.
The fundamental mathematical advantage of this basis is that it completely decouples the translation matrix $\mathbf{H}_{ij}$ between the two polarisation blocks.
While the single-sphere T-matrix is diagonal in the combined multipole index $mn = n(n+1) + m$ in both the standard TE/TM and lr\_tran bases, in the lr\_tran basis it couples the two polarisation blocks (labeled $p = 1$ for left and $p = 2$ for right).
This structure greatly simplifies the multi-sphere interaction equations, as the computationally expensive translation operations can be performed independently for each polarisation block.

%% ============================================================
\section{Single-Sphere Scattering: Lorenz--Mie Theory}
\label{sec:mie}

For a homogeneous sphere with size parameter $x$ and relative refractive index $m_\mathrm{rel}$, the Lorenz--Mie scattering coefficients $a_n$ (TM/electric) and $b_n$ (TE/magnetic) are given by (BH83, Eq.~4.88):
\begin{equation}
  a_n = \frac{\left[\dfrac{D_n(m_\mathrm{rel}\,x)}{m_\mathrm{rel}} + \dfrac{n}{x}\right] \psi_n(x) - \psi_{n-1}(x)}
             {\left[\dfrac{D_n(m_\mathrm{rel}\,x)}{m_\mathrm{rel}} + \dfrac{n}{x}\right] \xi_n(x) - \xi_{n-1}(x)}
  \label{eq:mie-an}
\end{equation}
\begin{equation}
  b_n = \frac{\left[m_\mathrm{rel}\,D_n(m_\mathrm{rel}\,x) + \dfrac{n}{x}\right] \psi_n(x) - \psi_{n-1}(x)}
             {\left[m_\mathrm{rel}\,D_n(m_\mathrm{rel}\,x) + \dfrac{n}{x}\right] \xi_n(x) - \xi_{n-1}(x)}
  \label{eq:mie-bn}
\end{equation}
where:
\begin{itemize}[nosep]
  \item $\psi_n(z) = z\,j_n(z)$ is the Riccati--Bessel function ($j_n$: spherical Bessel function of the first kind)
  \item $\xi_n(z) = z\,h_n^{(1)}(z) = \psi_n(z) + i\,\chi_n(z)$ is the Riccati--Hankel function
  \item $\chi_n(z) = -z\,y_n(z)$ ($y_n$: spherical Bessel function of the second kind)
  \item $D_n(z) = \frac{d}{dz}\ln[z\,j_n(z)] = \psi_n'(z)/\psi_n(z)$ is the logarithmic derivative
\end{itemize}

The logarithmic derivative $D_n(y)$ with $y = m_\mathrm{rel}\,x$ is computed by downward recurrence (Wiscombe 1980):
\begin{equation}
  D_{n-1}(y) = \frac{n}{y} - \frac{1}{D_n(y) + n/y},
  \label{eq:dn-recurrence}
\end{equation}
starting from a sufficiently large $n_\mathrm{start} \geq \max(N_\mathrm{max}, |y|) + 16$ where $D_{n_\mathrm{start}} \approx 0$.
This recurrence is unconditionally stable for any complex $y$.
The Riccati--Bessel functions $\psi_n(x)$ and $\chi_n(x)$ (for real argument $x$) are computed by upward recurrence:
\begin{equation}
  f_n(x) = \frac{2n - 1}{x}\,f_{n-1}(x) - f_{n-2}(x)
  \label{eq:bessel-recurrence}
\end{equation}
with starting values $\psi_0 = \sin x$, $\psi_1 = \sin x / x - \cos x$, $\chi_0 = -\cos x$, $\chi_1 = -\cos x / x - \sin x$.

The multipole expansion is truncated at order $N_i$ determined by the criterion of Wiscombe (1980) as implemented in the MSTM code:
\begin{equation}
  N_i = \mathrm{round}\!\left(x_i + 4\,x_i^{1/3}\right) + 5,
  \label{eq:nstop}
\end{equation}
refined by a convergence check: starting from this upper bound, $N_i$ is reduced until the cumulative partial extinction efficiency $Q_\mathrm{ext}^{(N_i)}$ deviates from the total by more than $10^{-6}$.

In the lr\_tran basis, the single-sphere T-matrix for multipole order $n$ takes the $2 \times 2$ form coupling the $p = 1$ and $p = 2$ blocks (Mackowski 1996, derived from the \texttt{mieoa} routine in MSTM v4.0):
\begin{equation}
  \mathbf{T}_n = -\frac{1}{2}
  \begin{pmatrix} a_n + b_n & a_n - b_n \\ a_n - b_n & a_n + b_n \end{pmatrix}
  \label{eq:tmatrix-lr}
\end{equation}
This matrix is diagonal in the combined index $mn$ (it does not mix different $(m,n)$ pairs) but couples the two circular polarisation blocks.
For an isotropic sphere, $a_n \neq b_n$ in general, so the off-diagonal coupling is physically meaningful and arises from the difference in electric and magnetic multipole scattering.

%% ============================================================
\section{Multi-Sphere Interaction Equation}
\label{sec:multi-sphere}

\subsection{Coupled field equations}

The field that excites sphere $i$ consists of the external incident wave plus the scattered fields from all other spheres, translated to the coordinate origin of sphere $i$ using the VSWF translation addition theorem (Mackowski 1991, 1996; Xu 1996):
\begin{equation}
  \bm{p}_\mathrm{exc}^{(i)} = \bm{p}_\mathrm{inc}^{(i)} + \sum_{j \neq i} \mathbf{H}_{ij}\,\bm{a}^{(j)}
  \label{eq:exciting-field}
\end{equation}
where $\mathbf{H}_{ij}$ is the translation matrix that re-expands outgoing VSWFs centred at $\bm{r}_j$ into regular VSWFs centred at $\bm{r}_i$ (Hankel-to-Bessel type), and $\bm{a}^{(j)}$ is the vector of scattered-field expansion coefficients of sphere $j$.

The T-matrix relates the exciting field to the scattered field:
\begin{equation}
  \bm{a}^{(i)} = \mathbf{T}_i\,\bm{p}_\mathrm{exc}^{(i)}.
  \label{eq:tmatrix-relation}
\end{equation}

Substituting and rearranging into a single global linear system for $\bm{x} = (\bm{a}^{(1)}, \dots, \bm{a}^{(N_S)})^\top$ gives (Mackowski 1994, Eq.~6):
\begin{equation}
  (\mathbf{I} - \mathbf{T}\,\mathbf{A})\,\bm{x} = \mathbf{T}\,\bm{p}_\mathrm{inc}
  \label{eq:global-system}
\end{equation}
where $\mathbf{T} = \mathrm{diag}(\mathbf{T}_1, \dots, \mathbf{T}_{N_S})$ is the block-diagonal T-matrix and $\mathbf{A}$ is the off-diagonal translation operator with blocks $\mathbf{H}_{ij}$ ($i \neq j$) and zeros on the diagonal.
The dimension of this system is $\sum_i 2\,N_i(N_i+2)$ --- twice the multipole block size because each sphere carries both $p = 1$ and $p = 2$ polarisation blocks.

\subsection{Translation matrix via rotation--axial translation--rotation}

The translation matrix $\mathbf{H}_{ij}$ is computed via the factorisation into rotation, axial translation, and inverse rotation (Mackowski 1996, \S 3; Xu 1996, \S 4), which avoids the direct evaluation of Gaunt integrals and is more efficient:

\begin{enumerate}[nosep]
  \item \textbf{Rotation.}
  Compute the spherical coordinates $(d_{ij}, \theta_{ij}, \phi_{ij})$ of the translation vector $\bm{r}_i - \bm{r}_j$.
  The azimuthal phase factors $e^{im\phi_{ij}}$ and the normalised associated Legendre functions $P_n^m(\cos\theta_{ij})$ provide the rotation from the laboratory frame to a frame where the translation vector is along $z$.

  \item \textbf{Axial translation coefficients.}
  In the rotated frame where the two origins are separated along $z$ by distance $d_{ij}$, the translation coefficients factorise.
  The key quantity involves a sum over intermediate order $w$:
  \begin{equation}
    C_{mn,kl} = \sum_{w=w_\mathrm{min}}^{n+l} c_w(m,n,k,l) \; h_w^{(1)}(d_{ij}) \; P_w^{m-k}(\cos\theta_{ij})
    \label{eq:axial-translation}
  \end{equation}
  where $w_\mathrm{min} = \max(|n-l|, |m-k|)$ and $c_w(m,n,k,l)$ are products of vector coupling (Gaunt) coefficients.
  These coupling coefficients are related to Clebsch--Gordan coefficients and are evaluated using a combined downward--upward recurrence algorithm (Xu 1996, \S 3; Mackowski 1996, Appendix) that avoids explicit factorial computation and maintains numerical stability even at high multipole orders.

  \item \textbf{Polarisation coupling in the lr\_tran basis.}
  In this basis the translation matrix has a particularly simple structure: the even-$w$ and odd-$w$ contributions (relative to $n + l$) combine to give the two polarisation blocks independently.
  Specifically, denoting the even- and odd-$w$ partial sums as $A$ and $B$:
  \begin{equation}
    H_{kl,mn}^{(p=1)} = e^{i(m-k)\phi_{ij}} \left(A + iB\right), \qquad
    H_{kl,mn}^{(p=2)} = e^{i(m-k)\phi_{ij}} \left(A - iB\right).
    \label{eq:lr-translation}
  \end{equation}
  This decoupling is the fundamental advantage of the lr\_tran basis (Mackowski 1996, Eq.~18): the two circular polarisation modes interact with each other only through the T-matrix, not through translation.
\end{enumerate}

The coupling coefficients are precomputed once and cached as a 3D array \texttt{tran\_coef[mn, kl, w]} (Fortran routine \texttt{gentrancoefconstants}), while the distance-dependent Hankel functions and Legendre functions are evaluated per sphere pair.

\subsection{FFT-accelerated translation}

For large numbers of spheres ($N_S \gtrsim 100$), the direct pairwise evaluation of $\mathbf{A}\,\bm{v}$ scales as $O(N_S^2)$ per iteration and dominates the computation time.
Mackowski \& Kolokolova (2022) introduced a discrete Fourier convolution (DFC) scheme into MSTM v4.0 that reduces this to approximately $O(N_S \ln N_S)$ scaling.
The algorithm proceeds as follows:

\begin{enumerate}[nosep]
  \item \textbf{Grid construction.}
  A uniform cubic grid with cell spacing $d$ is mapped onto the target volume.
  The cell spacing is chosen from the mean monomer radius and a target volume fraction $f_v$ as $d = \bar{a} \left(4\pi / 3 f_v\right)^{1/3}$, and each grid dimension is rounded up to the nearest integer of the form $2^a \times 3^b \times 5^c$ for efficient FFT computation (Temperton's generalised prime-factor FFT).

  \item \textbf{Sphere--node association.}
  Each sphere is assigned to the grid cell containing its centre.
  A \emph{neighbour hole} region is defined around each cell: for all cell pairs within this exclusion region (typically $|i - j| \leq \sqrt{2}$, i.e., self + face-adjacent + edge-adjacent cells), sphere--sphere interactions are computed by direct Hankel-type translation, preserving accuracy at short inter-sphere distances.

  \item \textbf{Sphere $\to$ Node.}
  For each sphere, the scattered-field expansion coefficients are translated to its associated grid node using a regular (Bessel $j_n$) translation matrix $\mathbf{J}$.
  This is valid because the sphere centre is within the grid cell (distance less than $d/2$), well within the convergence region of the regular addition theorem.

  \item \textbf{Node-to-node convolution.}
  The grid-based translation operator is a discrete convolution: the translation matrix from node $\bm{i}$ to node $\bm{j}$ depends only on the displacement $\bm{j} - \bm{i}$.
  This convolution is evaluated in the frequency domain:
  \begin{equation}
    \hat{g}_\mu(\bm{k}) = \sum_\nu \hat{H}_{\mu\nu}(\bm{k})\,\hat{a}_\nu(\bm{k})
    \label{eq:fft-convolution}
  \end{equation}
  where $\mu$ and $\nu$ represent the combined multipole indices.
  $\hat{H}_{\mu\nu}(\bm{k})$ is the pre-computed 3D FFT of the cell-to-cell translation matrix (excluding the neighbour-hole cells), $\hat{a}_\nu(\bm{k})$ is the FFT of the zero-padded node coefficient field, and the product is taken element-wise per index pair $(\mu, \nu)$.
  The use of a doubled grid ($2N_x \times 2N_y \times 2N_z$) with zero-padding avoids circular convolution artefacts.

  \item \textbf{Node $\to$ Sphere.}
  The resulting node-level exciting field is translated back to each sphere position via the reverse regular translation.

  \item \textbf{Near-field correction.}
  Direct sphere-to-sphere Hankel-type translation is applied for all sphere pairs within the neighbour-hole region (step 2), adding these contributions to the FFT-based far-field result.
\end{enumerate}

The total cost per iteration becomes $O(N_S + M \log M)$ where $M$ is the number of active grid cells, compared with $O(N_S^2 N_\mathrm{max}^4)$ for the direct algorithm.
The memory overhead consists of the pre-FFT'd translation matrices, which scale as $O(M \times N_\mathrm{node}^2)$ where $N_\mathrm{node}$ is the node multipole order.
Mackowski \& Kolokolova (2022) report speed improvements exceeding two orders of magnitude for systems of $10^3$--$10^4$ spheres.

%% ============================================================
\section{Iterative Solution: Complex BiConjugate Gradient}
\label{sec:cbicg}

The linear system $\mathbf{L}\,\bm{x} = \bm{b}$ where $\mathbf{L} = \mathbf{I} - \mathbf{T}\,\mathbf{A}$ and $\bm{b} = \mathbf{T}\,\bm{p}_\mathrm{inc}$ is solved by the Complex BiConjugate Gradient (CBICG) method (Mackowski \& Mishchenko 2011, \S 2.3).

\subsection{Operator definitions}

The CBICG iteration requires both the forward operator $\mathbf{L}\,\bm{v} = \bm{v} - \mathbf{T}\,\mathbf{A}\,\bm{v}$ and the adjoint (conjugate transpose) operator $\mathbf{L}^\dagger\,\bm{v} = \bm{v} - \mathbf{A}^\dagger\,\mathbf{T}^\dagger\,\bm{v}$.
The forward operator is computed by applying the translation $\mathbf{A}$ (direct or FFT-accelerated), followed by the block-diagonal T-matrix application $\mathbf{T}$.

The adjoint of the translation operator $\mathbf{A}$ is the computationally expensive part.
Mackowski \& Mishchenko (2011) exploit the S-parity symmetry of the VSWF translation matrix to avoid explicitly forming $\mathbf{A}^\dagger$:
\begin{equation}
  \mathbf{A}^\dagger = \left(\mathbf{S}\,\mathbf{A}\,\mathbf{S}\right)^*
  \label{eq:adjoint-symmetry}
\end{equation}
where $\mathbf{S}$ is the parity operator acting on each sphere's coefficient block as $a_{m,n,p} \mapsto (-1)^m\,a_{-m,n,p}$, and ${}^*$ denotes element-wise complex conjugation.
This identity arises from the spatial inversion symmetry and reciprocity of the electromagnetic translation operator.
The adjoint translation is evaluated by: (i)~applying $\mathbf{S}$, (ii)~performing the same forward translation $\mathbf{A}$, (iii)~applying $\mathbf{S}$ again, and (iv)~complex-conjugating.
This reuses the forward translation code path entirely, including the FFT-accelerated version, at no additional implementation cost.

The system is solved independently for two orthogonal incident linear polarisations ($q = 1$ corresponding to $x$-polarisation, and $q = 2$ corresponding to $y$-polarisation).
The default convergence tolerance is $\varepsilon = 10^{-6}$.

\subsection{CBICG algorithm (default solver)}

The CBICG (Complex BiConjugate Gradient) algorithm is summarised in Table~\ref{tab:cbicg}.  Each iteration requires one forward operator application $\mathbf{L}\,\bm{v}$ and one adjoint application $\mathbf{L}^\dagger\,\bm{v}$, totalling 2 translation operator ($\mathbf{A}$) evaluations per iteration.

\begin{table}[htbp]
\caption{CBICG algorithm (Mackowski \& Mishchenko 2011).}
\label{tab:cbicg}
\centering
\renewcommand{\arraystretch}{1.3}
\begin{tabular}{@{}p{0.95\textwidth}@{}}
\toprule
\textbf{Input:} $\mathbf{L}$, $\bm{b} = \mathbf{T}\,\bm{p}_\mathrm{inc}$, tolerance $\varepsilon$, max iterations $n_\mathrm{max}$. \\
\textbf{Initialise:} $\bm{x}_0 = \bm{b}$, \; $\bm{r}_0 = \bm{b} - \mathbf{L}\,\bm{x}_0$, \; $\bm{q}_0 = \overline{\bm{r}_0}$, \; $\bm{p}_0 = \bm{r}_0$, \; $\bm{w}_0 = \bm{q}_0$, \; $s_0 = \bm{q}_0^\mathsf{H} \bm{r}_0$. \\[4pt]
\textbf{For} $n = 0, 1, \ldots, n_\mathrm{max} - 1$: \\[2pt]
\quad 1. \; $\bm{u} = \mathbf{L}\,\bm{p}_n$ \hfill (1 forward operator application) \\
\quad 2. \; $\bm{v} = \mathbf{L}^\dagger\,\bm{w}_n$ \hfill (1 adjoint operator application via S-parity) \\
\quad 3. \; $\alpha_n = s_n \;/\; \bm{w}_n^\mathsf{H} \bm{u}$ \\
\quad 4. \; $\bm{x}_{n+1} = \bm{x}_n + \alpha_n\,\bm{p}_n$ \\
\quad 5. \; $\bm{r}_{n+1} = \bm{r}_n - \alpha_n\,\bm{u}$, \quad $\bm{q}_{n+1} = \bm{q}_n - \overline{\alpha}_n\,\bm{v}$ \\
\quad 6. \; $s_{n+1} = \bm{q}_{n+1}^\mathsf{H} \bm{r}_{n+1}$ \\
\quad 7. \; If $\mathrm{Re}(\bm{r}_{n+1}^\top \bm{r}_{n+1}) / \|\bm{b}\|^2 < \varepsilon$: \textbf{stop}, return $\bm{x}_{n+1}$. \\
\quad 8. \; $\beta_n = s_{n+1} / s_n$ \\
\quad 9. \; $\bm{p}_{n+1} = \bm{r}_{n+1} + \beta_n\,\bm{p}_n$, \quad $\bm{w}_{n+1} = \bm{q}_{n+1} + \overline{\beta}_n\,\bm{w}_n$ \\
\bottomrule
\end{tabular}
\end{table}

Note: the convergence criterion (step~7) uses the \emph{bilinear} (not sesquilinear) inner product $\bm{r}^\top \bm{r} = \sum_i r_i^2$, following the Fortran MSTM implementation.  The shadow residual $\bm{q} = \overline{\bm{r}}$ ensures that $\bm{q}^\mathsf{H}\bm{r} = \sum_i r_i^2$.

\subsection{GMRES algorithm (alternative solver)}

As an alternative to CBICG, a GMRES (Generalized Minimal Residual) solver (Saad \& Schultz 1986) is available via Krylov.jl (\texttt{solver=:gmres}).  GMRES requires only the forward operator $\mathbf{L}\,\bm{v}$ (no adjoint) and minimises the 2-norm of the residual over a Krylov subspace
\begin{equation}
  \mathcal{K}_k(\mathbf{L}, \bm{r}_0) = \mathrm{span}\{\bm{r}_0,\, \mathbf{L}\,\bm{r}_0,\, \mathbf{L}^2\,\bm{r}_0,\, \ldots,\, \mathbf{L}^{k-1}\,\bm{r}_0\}.
  \label{eq:krylov-subspace}
\end{equation}
The algorithm is summarised in Table~\ref{tab:gmres}.

\begin{table}[htbp]
\caption{GMRES algorithm (Saad \& Schultz 1986).}
\label{tab:gmres}
\centering
\renewcommand{\arraystretch}{1.3}
\begin{tabular}{@{}p{0.95\textwidth}@{}}
\toprule
\textbf{Input:} $\mathbf{L}$, $\bm{b}$, tolerance $\varepsilon$, restart parameter $m$, max iterations $n_\mathrm{max}$. \\
\textbf{Initialise:} $\bm{x}_0 = \bm{b}$, \; $\bm{r}_0 = \bm{b} - \mathbf{L}\,\bm{x}_0$, \; $\beta = \|\bm{r}_0\|_2$, \; $\bm{v}_1 = \bm{r}_0 / \beta$. \\[4pt]
\textbf{For} $j = 1, 2, \ldots, \min(m, n_\mathrm{max})$: \\[2pt]
\quad 1. \; $\bm{w} = \mathbf{L}\,\bm{v}_j$ \hfill (1 forward operator application) \\
\quad 2. \; \textbf{For} $i = 1, \ldots, j$: \; $h_{ij} = \bm{v}_i^\mathsf{H}\,\bm{w}$, \; $\bm{w} \leftarrow \bm{w} - h_{ij}\,\bm{v}_i$ \hfill (modified Gram--Schmidt) \\
\quad 3. \; $h_{j+1,j} = \|\bm{w}\|_2$, \; $\bm{v}_{j+1} = \bm{w} / h_{j+1,j}$ \\
\quad 4. \; Solve $\bm{y}_j = \arg\min_{\bm{y}} \|\beta\,\bm{e}_1 - \bar{\mathbf{H}}_j\,\bm{y}\|_2$ via Givens rotations \\
\quad 5. \; If $\|\beta\,\bm{e}_1 - \bar{\mathbf{H}}_j\,\bm{y}_j\|_2 / \|\bm{b}\|_2 < \varepsilon$: \textbf{stop}. \\[4pt]
\textbf{Solution update:} $\bm{x} = \bm{x}_0 + \mathbf{V}_j\,\bm{y}_j$, where $\mathbf{V}_j = [\bm{v}_1, \ldots, \bm{v}_j]$. \\
If $j = m$ and not converged: restart with $\bm{x}_0 \leftarrow \bm{x}$. \\
\bottomrule
\end{tabular}
\end{table}

Key properties of GMRES compared to CBICG:
\begin{itemize}
  \item 1 operator application per iteration (no adjoint), versus 2 for CBICG.
  \item Monotonically decreasing residual: $\|\bm{r}_k\|_2 \leq \|\bm{r}_{k-1}\|_2$ by construction.
  \item Memory: $O(kn)$ for $k$ Krylov basis vectors of length $n$ (${\sim}\SI{14}{\mega\byte}$ for $k = 30$, $n = 48000$).
  \item The implementation uses $m = 30$ (Krylov.jl \texttt{memory} parameter).
\end{itemize}

\textbf{Performance comparison.}
For the MSTM problem class, GMRES requires approximately twice as many iterations as CBICG to reach the same residual tolerance, because the CBICG adjoint operator (via S-parity, \eqref{eq:adjoint-symmetry}) provides additional directional information.  The total number of $\mathbf{A}$-operator applications is therefore similar (CBICG: $k$ iterations $\times$ 2; GMRES: $\sim 2k$ iterations $\times$ 1), and the wall-clock performance is comparable.  The default solver remains CBICG, consistent with the Fortran MSTM reference implementation.

\subsection{Continuation method for refractive index sweeps}
\label{sec:continuation}

When sweeping the complex refractive index $m$ over a grid while keeping the aggregate geometry fixed, the solution $\bm{x}(m)$ of \eqref{eq:global-system} varies continuously with $m$.  A natural acceleration strategy is to use the converged solution $\bm{x}(m_j)$ at grid point $m_j$ as the initial guess $\bm{x}_0$ for the iterative solver at the neighbouring grid point $m_{j+1}$, instead of the default initial guess $\bm{x}_0 = \bm{b} = \mathbf{T}\,\bm{p}_\mathrm{inc}$.  This is known as a \emph{continuation method} (or warm-start).

When the refractive index changes from $m_j$ to $m_{j+1}$, only the Mie coefficients $a_n, b_n$ (and hence the T-matrix $\mathbf{T}$) change; the translation operator $\mathbf{A}$ depends only on geometry and is unchanged.  For small $|m_{j+1} - m_j|$, the new solution is close to the previous one, so the initial residual $\|\bm{b}_{j+1} - \mathbf{L}_{j+1}\,\bm{x}(m_j)\|$ is much smaller than $\|\bm{b}_{j+1} - \mathbf{L}_{j+1}\,\bm{b}_{j+1}\|$, leading to fewer iterations.

The continuation is applied automatically in the parameter sweep (\texttt{run\_parameter\_sweep}): within each (aggregate, medium condition) group, the solver coefficients $\bm{x}$ from the previous refractive index grid point are passed as the initial guess to the next.  If the equation system size changes (e.g., due to a change in the converged Mie truncation order), the continuation is skipped and the default initial guess is used.

Benchmarks (1000 spheres, FFT mode, $\Delta m \approx 0.05$) show a reduction from 15 to 9--11 iterations per grid point, corresponding to a wall-time speedup of approximately $1.2\times$.  For smaller aggregates ($N_p \sim 100$) where the cold-start solver already converges in 3--4 iterations, the improvement is marginal.

\subsection{Incident plane wave coefficients}

For $z$-axis incidence ($\beta = 0$, $\alpha = 0$), only the $m = \pm 1$ multipole components are non-zero due to the azimuthal symmetry of a circularly decomposed plane wave.
In the lr\_tran basis, the incident field coefficients for sphere $i$ are:
\begin{equation}
  p_{mn,p=1,q=1}^{(i)} = i^{\,n+1}\sqrt{\frac{2n+1}{2}}\;\delta_{m,+1}\;e^{iz_i}, \qquad
  p_{mn,p=2,q=1}^{(i)} = -i^{\,n+1}\sqrt{\frac{2n+1}{2}}\;\delta_{m,-1}\;e^{iz_i}
  \label{eq:pinc-q1}
\end{equation}
for the first incident polarisation ($q = 1$), and
\begin{equation}
  p_{mn,p=1,q=2}^{(i)} = -i^{\,n+2}\sqrt{\frac{2n+1}{2}}\;\delta_{m,+1}\;e^{iz_i}, \qquad
  p_{mn,p=2,q=2}^{(i)} = -i^{\,n+2}\sqrt{\frac{2n+1}{2}}\;\delta_{m,-1}\;e^{iz_i}
  \label{eq:pinc-q2}
\end{equation}
for $q = 2$.
The phase factor $e^{iz_i}$ accounts for the spatial position of sphere $i$ along the propagation direction.

%% ============================================================
\section{Output Quantities}
\label{sec:output}

\subsection{Translation to a common origin}

After solving for the per-sphere scattered-field coefficients $\bm{a}^{(i)}$, all coefficients are re-expanded about a common origin $\bm{r}_0$ (typically the aggregate centroid) using regular (Bessel $j_n$) translation:
\begin{equation}
  \bm{a}_0 = \sum_{i=1}^{N_S} \mathbf{J}_{0i}\,\bm{a}^{(i)}
  \label{eq:common-origin}
\end{equation}
where $\mathbf{J}_{0i}$ is the regular translation matrix from sphere $i$ to the common origin.
The translation order for each sphere is determined by a convergence test on the translation addition theorem (Fortran routine \texttt{tranordertest}), ensuring that the re-expansion captures the scattered field to a specified accuracy.

The common-origin lr\_tran coefficients are then transformed from the $(p=1, p=2)$ basis to mode-indexed form:
\begin{equation}
  a_{mn}^{(\text{mode}\,1)} = a_{mn}^{(p=1)} + a_{mn}^{(p=2)}, \qquad
  a_{mn}^{(\text{mode}\,2)} = a_{mn}^{(p=1)} - a_{mn}^{(p=2)}.
  \label{eq:mode-transform}
\end{equation}
These mode-1 and mode-2 coefficients correspond to the standard TE and TM VSWF expansion of the total scattered field about the common origin (Mackowski 1994, \S 3).

\subsection{Scattering amplitude matrix (BH83 convention)}

The far-field scattered wave is described by the $2 \times 2$ amplitude matrix (BH83, Eq.~3.12):
\begin{equation}
  \begin{pmatrix} E_\parallel^\mathrm{sca} \\ E_\perp^\mathrm{sca} \end{pmatrix}
  = \frac{e^{ikr}}{-ikr}
  \begin{pmatrix} S_2 & S_3 \\ S_4 & S_1 \end{pmatrix}
  \begin{pmatrix} E_\parallel^\mathrm{inc} \\ E_\perp^\mathrm{inc} \end{pmatrix}
  \label{eq:bh83-amplitude}
\end{equation}
where $S_1, S_2, S_3, S_4$ are dimensionless complex functions of the scattering angles $(\theta, \phi)$.

For the forward ($\theta = 0^\circ$) and backward ($\theta = 180^\circ$) directions only the $m = \pm 1$ terms of the common-origin expansion survive.
The amplitudes are computed from the mode coefficients as
\begin{equation}
  S_j = -2 \sum_{n=1}^{N_\mathrm{max}} \sum_{m = \pm 1} \sum_{p = 1}^{2}
  F_j(n, m, p)\;\tau_{m,n,p}(\theta)
  \label{eq:amplitude-sum}
\end{equation}
where $F_j$ selects the appropriate combination of mode coefficients and incident polarisation (mode-1 from $q=1$, mode-2 from $q=1$ for $S_2$, $S_4$; similarly with sign change from $q=2$ for $S_1$, $S_3$), and $\tau_{m,n,p}(\theta)$ are angular functions derived from the Wigner $d$-matrix at $\cos\theta$.

At $\theta = 0^\circ$ the angular functions simplify to (Mackowski \& Mishchenko 2011, \texttt{taufunc} routine):

\begin{center}
\begin{tabular}{@{} ccc @{}}
  \toprule
  $m$ & $p$ & $\tau_{m,n,p}(0^\circ)$ \\
  \midrule
  $+1$ & 1 & $-f_n$ \\
  $+1$ & 2 & $-f_n$ \\
  $-1$ & 1 & $+f_n$ \\
  $-1$ & 2 & $-f_n$ \\
  \bottomrule
\end{tabular}
\end{center}
where $f_n = \frac{1}{4}\sqrt{(2n+1)/2}$.

At $\theta = 180^\circ$ an additional factor of $(-1)^n$ appears, and the signs of the $p = 1$ vs $p = 2$ contributions differ:

\begin{center}
\begin{tabular}{@{} ccc @{}}
  \toprule
  $m$ & $p$ & $\tau_{m,n,p}(180^\circ)$ \\
  \midrule
  $+1$ & 1 & $-(-1)^n f_n$ \\
  $+1$ & 2 & $+(-1)^n f_n$ \\
  $-1$ & 1 & $+(-1)^n f_n$ \\
  $-1$ & 2 & $+(-1)^n f_n$ \\
  \bottomrule
\end{tabular}
\end{center}

\textbf{Note on $S_3$ and $S_4$:}
For a single sphere or a symmetric dimer with the symmetry axis along $z$, $S_3 = S_4 = 0$ by symmetry.
For general aggregates lacking such symmetry, all four elements are non-zero.

\subsection{Cross-section efficiencies}

The dimensionless efficiency factors are defined as $Q = C / (\pi\,a_\mathrm{eff}^2)$ where $C$ is the corresponding cross section, and computed for unpolarised incidence by averaging over the two incident polarisations $q = 1, 2$.

\textbf{Extinction efficiency} (optical theorem, Mackowski 1994, Eq.~12):
\begin{equation}
  Q_\mathrm{ext} = -\frac{2}{x_\mathrm{eff}^2}
  \sum_{q=1}^{2} \mathrm{Re}\!\left(\bm{a}_q^\dagger \,\bm{p}_{\mathrm{inc},q}\right)
  \label{eq:qext}
\end{equation}
where the inner product $\bm{a}_q^\dagger \bm{p}_{\mathrm{inc},q}$ is evaluated in the per-sphere basis (before common-origin translation).
This is the multi-sphere generalisation of the forward-scattering (optical) theorem: the extinction is determined entirely by the projection of the scattered-field coefficients onto the incident-field pattern.
The factor of 2 (instead of 4) in the prefactor arises from averaging the cross section over the two incident orthogonal polarisation states ($q=1, 2$).

\textbf{Scattering efficiency} (from the coherent common-origin expansion, Mackowski 1994, Eq.~13):
\begin{equation}
  Q_\mathrm{sca} = \frac{1}{x_\mathrm{eff}^2}
  \sum_{q=1}^{2} \left(\|\bm{a}_{0,q}^{(\text{mode}\,1)}\|^2 + \|\bm{a}_{0,q}^{(\text{mode}\,2)}\|^2 \right)
  \label{eq:qsca}
\end{equation}
where the norm is the sum of squared magnitudes of all common-origin expansion coefficients for incident polarisation $q$.
This requires the common-origin re-expansion (Sec.~\ref{sec:output}, first subsection), which coherently accounts for the relative phases of all spheres.

\textbf{Absorption efficiency:}
\begin{equation}
  Q_\mathrm{abs} = Q_\mathrm{ext} - Q_\mathrm{sca}
  \label{eq:qabs}
\end{equation}
following from energy conservation.

\subsection{MI02 scattering amplitudes}

The MI02 (Mishchenko, Travis \& Lacis 2002) amplitude matrix $\mathbf{S}^\mathrm{MI02}$ has dimension of length and is related to the dimensionless BH83 amplitudes by:
\begin{equation}
  S_{11}^\mathrm{MI02} = \frac{S_2^\mathrm{BH83}}{-ik}, \quad
  S_{22}^\mathrm{MI02} = \frac{S_1^\mathrm{BH83}}{-ik}, \quad
  S_{12}^\mathrm{MI02} = \frac{S_3^\mathrm{BH83}}{ik}, \quad
  S_{21}^\mathrm{MI02} = \frac{S_4^\mathrm{BH83}}{ik}
  \label{eq:mi02-conversion}
\end{equation}
where $k = 2\pi\,n_\mathrm{med}/\lambda_0$ is the (dimensional) wavenumber in the medium.
Note the index mapping: the MI02 $S_{11}$ element corresponds to BH83 $S_2$ (not $S_1$), and vice versa.

\subsection{CAS-v2 observable amplitudes}

The Complex Amplitude Sensing (CAS-v2) protocol (Moteki \& Adachi 2024) directly measures three complex amplitudes derived from the MI02 matrix elements:
\begin{align}
  S_s^\mathrm{fwd} &= S_{11}^\mathrm{MI02} + i\,S_{12}^\mathrm{MI02} \label{eq:cas-ss} \\
  S_p^\mathrm{fwd} &= S_{22}^\mathrm{MI02} - i\,S_{21}^\mathrm{MI02} \label{eq:cas-sp} \\
  S^\mathrm{bak}   &= \frac{-S_{11}^\mathrm{MI02,bwd} + S_{22}^\mathrm{MI02,bwd}}{\sqrt{2}} \label{eq:cas-sbak}
\end{align}
where $S_s^\mathrm{fwd}$ and $S_p^\mathrm{fwd}$ are the forward scattering amplitudes for $s$- and $p$-polarised incident light, respectively, and $S^\mathrm{bak}$ is the depolarisation-sensitive backward scattering amplitude.

%% ============================================================
\section{Computational Flow Summary}
\label{sec:flow}

\begin{center}
\begin{tabular}{@{} rlll @{}}
  \toprule
  Step & Operation & Description & Source file \\
  \midrule
  1 & Non-dimensionalise & $x_i = k \cdot a_i$, $\bm{r}_i \to k \cdot \bm{r}_i$, $m_\mathrm{rel} = m_\mathrm{sphere}/n_\mathrm{med}$ & \texttt{ParameterSweep.jl} \\
  2 & Mie coefficients & $a_n, b_n$ for each sphere & \texttt{MieCoefficients.jl} \\
  3 & Build T-matrix & lr\_tran basis, per sphere & \texttt{TMatrixSolver.jl} \\
  4 & Incident field & $\bm{p}_\mathrm{inc}$ ($m = \pm 1$ terms) & \texttt{TMatrixSolver.jl} \\
  5 & CBICG solve & $(\mathbf{I} - \mathbf{T}\mathbf{A})\bm{x} = \mathbf{T}\bm{p}_\mathrm{inc}$; warm-start from previous $m$ (\S\ref{sec:continuation}) & \texttt{TMatrixSolver.jl} \\
  6 & Common origin & Translate coefficients & \texttt{ScatteringAmplitude.jl} \\
  7 & Amplitudes & $S_1$--$S_4$ at $\theta = 0^\circ, 180^\circ$ & \texttt{ScatteringAmplitude.jl} \\
  8 & Efficiencies & $Q_\mathrm{ext}$, $Q_\mathrm{sca}$, $Q_\mathrm{abs}$ & \texttt{ScatteringAmplitude.jl} \\
  9 & MI02 \& CAS & Convert amplitudes & \texttt{ParameterSweep.jl} \\
  \bottomrule
\end{tabular}
\end{center}

%% ============================================================
\section{Implementation Mapping}
\label{sec:mapping}

\begin{center}
\begin{tabular}{@{} p{4.2cm} p{2.8cm} p{4.2cm} p{3.5cm} @{}}
  \toprule
  Theory & Equation(s) & Implementation & File \\
  \midrule
  Wavenumber, non-dim. & \eqref{eq:wavenumber}--\eqref{eq:a-eff} & \texttt{compute\_scattering}, \texttt{run\_parameter\_sweep} & \texttt{ParameterSweep.jl} \\
  Mie coefficients $a_n, b_n$ & \eqref{eq:mie-an}--\eqref{eq:nstop} & \texttt{mie\_coefficients} & \texttt{MieCoefficients.jl} \\
  lr\_tran T-matrix & \eqref{eq:tmatrix-lr} & \texttt{build\_tmatrix} & \texttt{TMatrixSolver.jl} \\
  Translation $\mathbf{H}_{ij}$ & \eqref{eq:axial-translation}--\eqref{eq:lr-translation} & \texttt{translate\_coefs} & \texttt{TranslationCoefs.jl} \\
  FFT translation & \eqref{eq:fft-convolution} & \texttt{fft\_translate}, \texttt{setup\_fft\_grid} & \texttt{FFTTranslation.jl} \\
  CBICG solver & \eqref{eq:global-system}, \eqref{eq:adjoint-symmetry}, Table~\ref{tab:cbicg} & \texttt{\_solve\_bicg} & \texttt{TMatrixSolver.jl} \\
  GMRES solver (optional) & \eqref{eq:global-system}, \eqref{eq:krylov-subspace}, Table~\ref{tab:gmres} & \texttt{gmres} (Krylov.jl) & \texttt{TMatrixSolver.jl} \\
  RI continuation (\S\ref{sec:continuation}) & --- & \texttt{initial\_amn} kwarg & \texttt{TMatrixSolver.jl}, \texttt{ParameterSweep.jl} \\
  Incident coefficients & \eqref{eq:pinc-q1}--\eqref{eq:pinc-q2} & \texttt{incident\_coefficients} & \texttt{TMatrixSolver.jl} \\
  Common-origin expansion & \eqref{eq:common-origin}--\eqref{eq:mode-transform} & \texttt{translate\_to\_common\_origin} & \texttt{ScatteringAmplitude.jl} \\
  BH83 amplitudes $S_1$--$S_4$ & \eqref{eq:bh83-amplitude}--\eqref{eq:amplitude-sum} & \texttt{compute\_amplitude\_matrix} & \texttt{ScatteringAmplitude.jl} \\
  Cross-section efficiencies & \eqref{eq:qext}--\eqref{eq:qabs} & \texttt{compute\_efficiencies} & \texttt{ScatteringAmplitude.jl} \\
  MI02 conversion & \eqref{eq:mi02-conversion} & \texttt{bh83\_to\_mi02} & \texttt{ParameterSweep.jl} \\
  CAS-v2 observables & \eqref{eq:cas-ss}--\eqref{eq:cas-sbak} & \texttt{read\_results\_h5\_for\_cas.jl} & \texttt{scripts/} \\
  Aggregate I/O & --- & \texttt{read\_aggregate\_file}, \texttt{read\_aggregate\_catalog} & \texttt{AggregateIO.jl} \\
  \bottomrule
\end{tabular}
\end{center}

%% ============================================================
\section{References}
\label{sec:references}

\begin{enumerate}[nosep]
  \item Bohren, C.F.\ \& Huffman, D.R.\ (1983). \textit{Absorption and Scattering of Light by Small Particles}. Wiley. [BH83]
  \item Mackowski, D.W.\ (1991). Analysis of radiative scattering for multiple sphere configurations. \textit{Proc.\ R.\ Soc.\ Lond.\ A}, 433, 599--614.
  \item Mackowski, D.W.\ (1994). Calculation of total cross sections of multiple-sphere clusters. \textit{JOSA A}, 11(11), 2851--2861.
  \item Mackowski, D.W.\ (1996). Calculation of the total scattering cross section and the total absorbed energy for the multiple sphere clusters. \textit{JOSA A}, 13(11), 2266--2278.
  \item Mackowski, D.W.\ (2014). A general superposition solution for electromagnetic scattering by multiple spherical domains of optically active media. \textit{JQSRT}, 133, 264--270.
  \item Mackowski, D.W.\ \& Kolokolova, L.\ (2022). Application of the multiple sphere superposition solution to large-scale systems of spheres via an accelerated algorithm. \textit{JQSRT}, 280, 108075.
  \item Mackowski, D.W.\ \& Mishchenko, M.I.\ (2011). A multiple sphere T-matrix Fortran code for use on parallel computer clusters. \textit{JQSRT}, 112(13), 2182--2192.
  \item Mishchenko, M.I., Travis, L.D.\ \& Lacis, A.A.\ (2002). \textit{Scattering, Absorption, and Emission of Light by Small Particles}. Cambridge University Press. [MI02]
  \item Moteki, N.\ \& Adachi, K.\ (2024). CAS-v2 protocol. \textit{Optics Express}, 32(21), 36500--36522.
  \item Saad, Y.\ \& Schultz, M.H.\ (1986). GMRES: A generalized minimal residual algorithm for solving nonsymmetric linear systems. \textit{SIAM J.\ Sci.\ Stat.\ Comput.}, 7(3), 856--869.
  \item Wiscombe, W.J.\ (1980). Improved Mie scattering algorithms. \textit{Applied Optics}, 19(9), 1505--1509.
  \item Xu, Y.-l.\ (1996). Calculation of the addition coefficients in electromagnetic multisphere-scattering theory. \textit{J.\ Comput.\ Phys.}, 127(2), 285--298.
\end{enumerate}

\bigskip
\noindent\textbf{Acknowledgment.}\quad
This document was prepared with the assistance of Claude (Anthropic). The author assumes full responsibility for the content.

\end{document}
